\chapter{Arhitectura Sistemului}
\label{ch:arhitectura}

În procesul dezvoltării sistemului software pentru proiectul \textit{Farm Defense: Harvest Wars}, a fost esențială utilizarea unor șabloane de proiectare (design patterns) consacrare și a
unei arhitecturi bine structurate. Arhitectura aleasă asigură separarea responsabilităților și mentenabilitatea pe termen lung a codului.

\section{Arhitectura Generală a Sistemului (Modelul C4)}
\label{sec:arhitectura_sistem}

Pentru a vizualiza arhitectura sistemului distribuit \textit{Farm Defense: Harvest Wars}, a fost utilizat modelul C4 \cite{brown_c4_model} la nivelul 2, diagramă de containere (Container Level). Modelul C4 a fost creat de Simon Brown ca o alternativă modernă la diagrama clasică UML. Avantajele sale includ simplitatea și claritatea reprezentării arhitecturii software. În acest caz, sistemul este structurat precum un model client-server, unde exista decuplare clară între logica de business și faza de gameplay în timp real.

În figura \ref{fig:c4_container}, se observă arhitectura formată din 4 containere principale:
\begin{enumerate}
    \item \textbf{Godot Game Client:} Reprezintă aplicația ce există pe computerul utilizatorului. Prin intermediul ei ultilizatorul interacționează atât cu logica din meniu (ex: selectarea unităților, deschidera cutiilor de recompense etc.), cât și cu faza de gameplay în timp real (ex: plasarea unităților etc.)
    \item \textbf{ASP.NET Core Backend:} Funcționează ca un REST API și ca un orchestrator de instanțe de joc. Acesta validează autentificarea utilizatorilor prin sistemul Identity, procesează tranzacțiile de tip metagame (modificarea pachetelor de unități) și rulează algoritmul de matchmaking.
    \item \textbf{Baza de Date SQLite:}
          Asigură stratul de persistență a datelor, comunicarea este abstractizată prin pachetul \textbf{Entity Framework Core}, care oferă un model de programare orientat pe obiecte.
    \item \textbf{Godot Headless Servers (Servere Dedicate):} Când se formează un meci valid, backend-ul execută un nou proces Godot rulat în mod \textit{headless} (fără interfață grafică), câruia îi sunt transmise datele despre meci si jucatori(id-ul meciului, unitătile selectate de jucatori etc.). Clientul primește datele meciului pentru a se putea conecta direct la acest server dedicat pentru faza de \textit{gameplay}.
\end{enumerate}

Această abordare micro-arhitecturală previne blocajele (bottlenecks) la nivelul API-ului web și izolează logic fiecare meci în propriul proces alocat de sistemul de operare Linux.

\begin{figure}[htbp]
    \centering
    \includegraphics[width=0.7\textwidth]{DiagramaC4.png}

    \caption{Diagrama de Arhitectură C4 (Nivel Container)}
    \label{fig:c4_container}
\end{figure}

\section{Arhitectura Clientului (Godot Engine)}
\label{sec:arhitectura_client}

Clientul este construit arhitectural pe principiul modularității. Principalul design pattern utilizat este \textit{Component Pattern}\cite{nystrom_game_patterns}, adaptat la specificul motorului de joc Godot. În loc să folosim clase monolitice pentru unitățile de joc, am optat pentru o compoziție de noduri copii specializate.


\subsection{Șablonul Component (Component Pattern)}

Fiecare unitate (ex: \textit{ChickenUnit}) este un ansamblu de componente independente precum \texttt{HitboxComponent}, \texttt{HealthComponent} și \texttt{UnitStateMachine}. Această abordare permite atât reutilizarea codului cât și testarea fiecărei funcționalități în mod independent.


\begin{figure}[htbp]
    \centering
    \includegraphics[width=0.9\textwidth]{ChickenC4Component.png}
    \caption{Diagrama de Componente a unei Unități, nivel Component \cite{brown_c4_model}}
    \label{fig:diagrama_componente}
\end{figure}

Pe baza arhitecturii ilustrate, responsabilitățile din entitatea \textit{ChickenUnit} au fost împărțite, conform principiului SRP (Single Responsibility Principle) \cite{martin2003agile}, în următoarele componente:

\begin{itemize}
    \item \textbf{\texttt{UnitStateMachine}:} Reprezintă componenta care orchestrează tranzițiile între stările de idle, move, attack și die.
    \item \textbf{\texttt{HealthComponent} și \texttt{HurtboxComponent}:} Gestionează variabila de viață a unei unități și validează coliziunile fizice ce au rol în aplicarea daunelor.
    \item \textbf{\texttt{VisionComponent} (Principal și Secundar):} Acționează ca senzori pentru a detecta unitățile inamice. Unitățile Defender au doi senzori de viziune: unul în partea stângă și unul în partea dreaptă.
    \item \textbf{\texttt{UnitVisualsComponent}:} Funcționează pe bază de evenimente (semnale), ascultând schimbările de stare și actualizează animațiile și HUD-ul (Health Bar).
\end{itemize}

\subsection{Șablonul Stare (State Pattern)}
\label{sec:state_pattern}

Comportamentul unităților de joc a fost modelat în principiu folosind șablonul de proiectare Stare (State Pattern) \cite{gof_design_patterns}. De regulă, în cadrul unui joc de tip strategie, o unitate oscilează constant între diverse stări: repaus, deplasare sau atac. Abordarea naivă pentru modelarea acestui tip de comportament presupune utilizarea unor structuri condiționale masive (ex. instrucțiuni \texttt{switch} sau \texttt{if-else} înlănțuite) în cadrul loop-ului principal de actualizare (\texttt{\_PhysicsProcess} în Godot). Această tip de abordare duce rapid la o degradare a gestionării codului, încălcând principiul de responsabilitate unică (Single Responsibility Principle).

Soluția este implementată cu ajutorul clasei \texttt{UnitStateMachine}, care permite decuplarea logicii fiecărui comportament. Starea este reprezentată de o singură clasă separată care implementează interfața \texttt{IState}.

În acest caz particular, sunt definite trei comportamente(\ref{fig:state_transition}) de bază pentru unități, fiecare comportament fiind reprezentat de o clasă specifică:

% --- PLACEHOLDER DIAGRAMA TRANZIȚIE STĂRI ---
\begin{figure}[htbp]
    \centering
    \includegraphics[width=1\textwidth]{StateDiagram.png}
    \caption{Diagrama de tranziție a stărilor pentru unități}
    \label{fig:state_transition}
\end{figure}
% --------------------------------------------

\begin{itemize}
    \item \textbf{IdleState:} Gestionează starea în care unitatea este în repaus și scanează mediul pentru detectarea țintelor cu ajutorul componentei de viziune.
    \item \textbf{MovingState:} Controlează deplasarea în direcția bazei inamice sau către o unitate inamică.
    \item \textbf{AttackState:} Reprezintă componenta de luptă, este declanșată când o unitate inamică intră în raza de viziune.
\end{itemize}

\subsubsection{Specializarea stărilor prin Interfețe Derivate (Polimorfism)}
\label{sec:state_inheritance}

Pentru a nu duplica logica și pentru a gestiona moduri diferite de atac (atacatori corp-la-corp vs. atacatori la distanță), s-a recurs la implementarea unei ierarhii de tipuri bazată pe interfețe.

S-a introdus interfața \texttt{IAttackState}, care acționează ca o interfață marcaj (marker interface) și extinde interfața de bază \texttt{IState}. Aceasta definește contractul strict pentru orice secvență de luptă, expunând timpii de pregătire (\texttt{WindUpTimer}), de reîncărcare (\texttt{CooldownTimer}) și metoda declanșatoare \texttt{StartAttackSequence()}.
Arhitectura claselor care implementează acest sistem este detaliată în Figura \ref{fig:uml_attack_state}.

% --- PLACEHOLDER DIAGRAMA UML CLASE ---
\begin{figure}[htbp]
    \centering
    \includegraphics[width=0.9\textwidth]{ClassDiagram.png}
    \caption{Diagrama UML de clase ilustrând moștenirea interfeței IAttackState}
    \label{fig:uml_attack_state}
\end{figure}
% --------------------------------------

Execuția propriu-zisă a atacului este delegată claselor concrete care implementează această interfață:
\begin{itemize}
    \item \textbf{\texttt{MeleeAttackState}:} Implementează \texttt{IAttackState} și aplică daune directe prin intermediul componentei \textit{Hitbox}.
    \item \textbf{\texttt{RangedAttackState}:} Implementează \texttt{IAttackState} și se ocupă de instanțierea proiectilelor.
\end{itemize}


\section{Arhitectura Backend-ului (.NET 10)}
\label{sec:arhitectura_backend}

Serverul este organizat folosind o arhitectură stratificată (Layered Architecture), asigurând o separare clară între interfața API, logica de business și accesul la date.

Pentru gestionarea utilizatorilor și asigurarea securității datelor (hash-uirea parolelor, generarea de token-uri), sistemul backend integrează framework-ul \textit{ASP.NET Core Identity}\cite{aspnet_identity_docs}. Această soluție a fost aleasă deoarece oferă nativ o arhitectură robustă pentru autentificare și stocare securizată, compatibilă direct cu Entity Framework Core .

\begin{itemize}
    \item \textbf{Controllere:} Gestionează cererile HTTP și returnează răspunsuri formatate JSON.
    \item \textbf{Servicii:} Conțin logica centrală, cum ar fi \texttt{MatchmakingService}, care este injectat în controllere folosind mecanismul nativ de \textit{Dependency Injection} (DI) din .NET.
    \item \textbf{Entități:} Reprezintă structura datelor în cod și sunt mapate peste tabelele bazei de date.
\end{itemize}

Utilizarea DI permite un cuplaj slab între componente, facilitând testarea unitară și mock-uirea bazei de date sau a altor servicii externe.

\section{Modelarea Datelor (Baza de Date SQLite)}
\label{sec:database}

Baza de date constituie stratul de persistență al arhitecturii. Am ales SQLite pentru simplitatea sa (serverless) și performanța adecvată pentru un proiect aflat în stadiul de prototip și dezvoltare locală.

Interacțiunea cu baza de date este realizată prin \textbf{Entity Framework Core}, utilizând abordarea \textit{Code-First}. Aceasta ne permite să definim schema bazei de date direct prin clase C\#, generând ulterior migrațiile necesare.

Principalele entități modelate sunt:
\begin{itemize}
    \item \texttt{AspNetUsers:} Extinde utilizatorul standard Identity pentru a include atribute specifice jocului, precum cantitatea de aur (\textit{Gold}), punctele de experiență (\textit{XP}) și nivelul.
    \item \texttt{Decks:} Reprezintă seturile de unități configurate de jucător pentru luptă, având o relație de tip 1-la-N cu utilizatorul.
    \item \texttt{UnitUnlocks:} Gestionează colecția de unități deblocate de fiecare jucător, permițând o progresie controlată în joc.
\end{itemize}

\begin{figure}[htbp]
    \centering
    % \includegraphics[width=0.9\textwidth]{erd_diagram}
    \caption{Diagrama Entitate-Relație (ERD)}
    \label{fig:erd_diagram}
\end{figure}

